% Parent document supplies amsmath, amssymb, amsthm, and theorem environments.
% Bibliography key ACEMKM is specified in gamma_provenance.txt.
\section{Reflected log-gamma moments and their late coefficients}
\label{reflected:sec:reflected-moments}

Set $f(x)=\log\Gamma(x)$ for $0<x<1$, and define
\[
 M_{n,m}=\int_0^1 f(x)^n f(1-x)^m\,dx,
 \qquad n,m\in\mathbb Z_{\ge0}.
\]
These integrals already occur as $T_{n+m,n}$ in
Amdeberhan et al.\ \cite[Section~5]{gaussian:ACEMKM}.  Their symmetry
$M_{n,m}=M_{m,n}$ and the identities obtained by expanding Euler's
reflection formula are therefore classical.  The asymptotic question here
is different: keep the reflected exponent $m$ fixed and let $n$ tend to
infinity.  The factor $f(1-x)^m$ vanishes to order $m$ at the endpoint
that produces the factorial growth.  This shifts the first exponential
scale from $1$ to $m+1$ without removing the eventual divergence of the
full expansion.  Theorem~\ref{reflected:gam:thm:late} identifies its late terms,
including their signs and the first relative correction.

\subsection{The meromorphic transform and all finite asymptotic orders}

Put
\[
 B(x)=\log\Gamma(1-x)
 =\gamma x+\sum_{j\ge2}\frac{\zeta(j)}j x^j,
 \qquad |x|<1,
\]
where $\gamma$ is Euler's constant.  For $k\ge1$ define
\begin{equation}\label{reflected:gam:eq:residue-coefficients}
 A_{k,m}=\frac1k[x^{k-1}]\exp\{kB(-x)\}B(x)^m.
\end{equation}
The convention $B(x)^0=1$ is in force.  In particular $A_{k,m}=0$
when $k\le m$.  Every coefficient is a finite polynomial over
$\mathbb Q$ in $\gamma,\zeta(2),\ldots,\zeta(k-1)$.

\begin{theorem}[Reflected moment expansion]\label{reflected:gam:thm:expansion}
For each fixed integer $m\ge0$, the function
\[
 F_m(z)=\int_0^1\Gamma(x)^z f(1-x)^m\,dx
\]
is holomorphic for $\operatorname{Re}z<m+1$, has Taylor expansion
\[
 F_m(z)=\sum_{n\ge0}\frac{M_{n,m}}{n!}z^n
 \qquad (|z|<m+1),
\]
and extends meromorphically to $\mathbb C$.  Its only possible poles
are simple poles at integers $k\ge m+1$, with
\begin{equation}\label{reflected:gam:eq:poles}
 \operatorname*{Res}_{z=k}F_m(z)=-kA_{k,m}.
\end{equation}
The first pole is present.  For every fixed integer $K\ge m+1$,
\begin{equation}\label{reflected:gam:eq:fixed-expansion}
 \frac{M_{n,m}}{n!}
 =\sum_{k=m+1}^{K}\frac{A_{k,m}}{k^n}
   +O_{m,K}\bigl((K+1)^{-n}\bigr),\qquad n\longrightarrow\infty.
\end{equation}
Consequently
\begin{equation}\label{reflected:gam:eq:leading}
 M_{n,m}\sim \frac{\gamma^m n!}{(m+1)^{n+1}}.
\end{equation}
\end{theorem}

\begin{proof}
As $x\downarrow0$, $f(x)=-\log x+O(x)$ and
$B(x)=\gamma x+O(x^2)$.  At the other endpoint,
$f(x)=O(1-x)$ and $f(1-x)=O(1+|\log(1-x)|)$.
These estimates prove absolute convergence and compact domination in
the asserted half-plane, including after any number of differentiations
in $z$.  They also justify the Taylor expansion near zero.

Fix $0<a<1$ and an integer $J\ge1$.  Taylor expansion in $x$ gives
\[
 \Gamma(1+x)^z B(x)^m
 =\sum_{j=0}^{J-1}b_{j,m}(z)x^j+x^J E_{J,m}(x,z),
\]
where each $b_{j,m}$ is an entire function of $z$, and the remainder is
uniformly bounded for $0\le x\le a$ and $z$ in compact sets.
The integral over $[0,a]$ is therefore
\[
 \sum_{j=0}^{J-1}\frac{b_{j,m}(z)a^{j+1-z}}{j+1-z}
 +\int_0^a x^{J-z}E_{J,m}(x,z)\,dx.
\]
The last integral is holomorphic for $\operatorname{Re}z<J+1$.
The integral over $[a,1]$ is entire, since its logarithmic singularity
at $1$ is integrable uniformly for $z$ in compact sets.  Letting $J$
increase gives the continuation.  At $z=k$ its residue is
$-b_{k-1,m}(k)=-kA_{k,m}$.  Since $B(x)^m$ starts with
$\gamma^m x^m$, the first nonzero coefficient is
$A_{m+1,m}=\gamma^m/(m+1)>0$.

For the asymptotic assertion, subtract the polar functions
\[
 \sum_{k=m+1}^{K+1}\frac{kA_{k,m}}{k-z}
\]
from $F_m(z)$.  The difference is holomorphic on $|z|<K+2$.
Cauchy's estimate on any circle $K+1<R<K+2$ bounds its $n$th Taylor
coefficient by $O_{m,K,R}(R^{-n})$.  Restoring the term at $K+1$
proves \eqref{reflected:gam:eq:fixed-expansion}, even if that particular coefficient
happens to vanish.  Taking $K=m+1$ proves \eqref{reflected:gam:eq:leading}.
\end{proof}

The first three coefficients are explicit.  Write $k=m+3$ in the last line:
\begin{align}
 A_{m+1,m}&=\frac{\gamma^m}{m+1},\label{reflected:gam:eq:first-three}\\
 A_{m+2,m}&=\frac{\gamma^{m-1}}{m+2}
 \left(\frac{m\zeta(2)}2-(m+2)\gamma^2\right),\notag\\
 A_{m+3,m}&=\frac{\gamma^m}{k}
 \left\{\frac{m\zeta(3)}{3\gamma}
 +\frac{m(m-1)\zeta(2)^2}{8\gamma^2}
 +\frac{k(1-m)\zeta(2)}2+\frac{k^2\gamma^2}2\right\}.\notag
\end{align}
These displays also apply to $m=0$, since $\gamma>0$ and factors of
$m$ remove the superfluous terms.  For example,
\begin{align*}
 \frac{M_{n,1}}{n!}
 &=\frac{\gamma}{2^{n+1}}
 +\frac{\zeta(2)/2-3\gamma^2}{3^{n+1}}
 +\frac{\zeta(3)/3+8\gamma^3}{4^{n+1}}+O(5^{-n}),\\
 \frac{M_{n,2}}{n!}
 &=\frac{\gamma^2}{3^{n+1}}
 +\frac{\gamma\zeta(2)-4\gamma^3}{4^{n+1}}+O(5^{-n}).
\end{align*}
The second displayed correction is positive, whereas the corresponding
correction for $m=1$ is negative.  Thus strict alternation from the first
term cannot be imported from the unreflected family.  The two signs can,
for instance, be checked using the elementary enclosures
$0.577<\gamma<0.578$ and $1.644<\zeta(2)<1.646$.
Eventual alternation has a different, precise rule.

\subsection{Late coefficients and a universal inverse-gamma scale}

Let $\phi(u)=\Gamma(1-u)$ and $D=u\,d/du$.  Define $\tau\in(0,1)$ by
\begin{equation}\label{reflected:gam:eq:critical}
 D\log\phi(\tau)=-\tau\psi(1-\tau)=1,
\end{equation}
and put
\begin{equation}\label{reflected:gam:eq:constants}
 \rho=\frac{\tau}{\phi(\tau)},\quad
 v=D^2\log\phi(\tau)=1+\tau^2\psi'(1-\tau),\quad
 C=\frac{\tau}{\sqrt{2\pi v}},\quad
 b=-\log\Gamma(1+\tau).
\end{equation}
The equation for $\tau$ has a unique solution: its left side equals
$\gamma u+\sum_{j\ge2}\zeta(j)u^j$, which is strictly increasing from
zero to infinity.  Also $0<\rho<1$ and $b>0$, the latter by strict
log-convexity of $\Gamma$ between $1$ and $2$.  Numerically,
$\tau=0.5040830082\ldots$ and $\rho=0.2821158090\ldots$.
Only the exact definitions are used below.

\begin{theorem}[Late reflected coefficients]\label{reflected:gam:thm:late}
For every fixed nonnegative integer $m$,
\begin{equation}\label{reflected:gam:eq:late}
 A_{k,m}=(-1)^{k+m-1}C b^m\rho^{-k}k^{-3/2}
 \left(1+\frac{\beta_m}{k}+O_m(k^{-2})\right).
\end{equation}
There is a full asymptotic expansion in integer powers of $1/k$.
To specify the first correction, set
$G(u)=\log\Gamma(1+u)$, $H_m(u)=G(u)^m$, and
$\kappa_j=D^j\log\phi(\tau)$.  All quantities on the right below
are evaluated at $u=\tau$:
\begin{equation}\label{reflected:gam:eq:beta}
 \beta_m=
 -\frac{(D+1)^2H_m}{2vH_m}
 +\frac{\kappa_3(D+1)H_m}{2v^2H_m}
 +\frac{\kappa_4}{8v^2}
 -\frac{5\kappa_3^2}{24v^3}.
\end{equation}
In particular, every sufficiently large possible pole in
Theorem~\ref{reflected:gam:thm:expansion} is present.  More explicitly,
\[
 \operatorname{Res}_{z=k}F_m(z)
 =(-1)^{k+m}C b^m\rho^{-k}k^{-1/2}
   \left(1+\frac{\beta_m}{k}+O_m(k^{-2})\right),
\]
so the eventual residue sign is $(-1)^{k+m}$.  For every fixed real $n$, the formal series
$\sum_{k\ge m+1}A_{k,m}/k^n$ diverges because its terms do not tend to zero.
\end{theorem}

\begin{proof}
Replacing $x$ by $-u$ in \eqref{reflected:gam:eq:residue-coefficients} gives
\[
 A_{k,m}=\frac{(-1)^{k-1}}k[u^{k-1}]\phi(u)^k H_m(u).
\]
The function $H_m$ is holomorphic for $|u|<1$.  If
$\phi(u)=\sum_{j\ge0}a_ju^j$, all the $a_j$ are strictly positive,
since $\log\phi$ has positive coefficients.  Thus
\[
 p_j=\frac{a_j\tau^j}{\phi(\tau)}
\]
is a probability distribution with exponential moments, mean one,
variance $v>0$, and positive mass at both $0$ and $1$.  It follows
that the function
\[
 Q(t)=\frac{\phi(\tau e^{it})}{\phi(\tau)}e^{-it}
\]
has modulus strictly less than one for $0<|t|\le\pi$, while
\[
 \log Q(t)=-\frac v2t^2-\frac{i\kappa_3}6t^3
 +\frac{\kappa_4}{24}t^4+O(t^5)
\]
near zero.  Cauchy's coefficient formula becomes
\begin{equation}\label{reflected:gam:eq:cauchy}
 A_{k,m}=
 \frac{(-1)^{k-1}\tau\rho^{-k}}{2\pi k}
 \int_{-\pi}^{\pi}Q(t)^k e^{it}H_m(\tau e^{it})\,dt.
\end{equation}

Fix a sufficiently small $\delta>0$.  On $\delta\le|t|\le\pi$
the integral is $O(q^k)$ for some $q<1$.  On
$k^{-2/5}\le|t|\le\delta$ it is $O(e^{-c k^{1/5}})$, because
$|Q(t)|\le e^{-ct^2}$.  The analytic amplitude is bounded on the
whole circle.  On the remaining interval set $t=y/\sqrt{k}$.
Gaussian domination and Taylor expansion to any fixed order, with
an analytic remainder, give a full expansion.  Odd terms integrate
to zero, so only integer powers of $1/k$ occur.  The leading integral is
\[
 H_m(\tau)\sqrt{\frac{2\pi}{vk}}
 =(-1)^m b^m\sqrt{\frac{2\pi}{vk}},
\]
which proves the leading term in \eqref{reflected:gam:eq:late}.

For the first correction, the amplitude derivatives at $t=0$ are
$H_m$, $i(D+1)H_m$, and $-(D+1)^2H_m$.
The Gaussian moments of degrees $2$, $4$, and $6$ contribute,
respectively, the amplitude's quadratic term, the amplitude--cubic
cross term together with the quartic phase term, and the square of
the cubic phase term.  Their normalized sum is exactly
\eqref{reflected:gam:eq:beta}.  Expanding two further even orders gives a
relative $O_m(k^{-2})$ remainder after the displayed correction.
Since $C b^m>0$ and $\rho<1$, the assertions about nonvanishing,
signs, and divergence follow at once.
\end{proof}

This proof uses the circle $|u|=\tau$ and a strict modulus inequality;
it does not presume a classification of all complex critical values
of the reciprocal gamma function.  The same circle governs every fixed
reflected exponent.  With $\alpha=-\log\rho$, the continuous minimum
of the leading term magnitude $\rho^{-k}k^{-n-3/2}$ occurs at
$k=(n+3/2)/\alpha$.  A least-term calculation alone does not bound the
remainder.  The next result supplies an explicit bound with a growing
number of terms.

\subsection{An effective remainder despite the initial sign changes}

\begin{theorem}[Explicit reflected truncation bound]\label{reflected:gam:thm:effective}
Let $m\ge0$, $n\ge1$, and $K\ge1$ be integers, put
$\alpha=-\log\rho$ and $T_m=\tau B(\tau)^m$, and suppose $K\alpha<n$.
With
\[
 R^{(m)}_{n,K}=\frac{M_{n,m}}{n!}
             -\sum_{k=1}^K\frac{A_{k,m}}{k^n},
 \qquad \Gamma(n,s)=\int_s^\infty t^{n-1}e^{-t}\,dt,
\]
one has the entirely explicit estimate
\begin{align}\label{reflected:gam:eq:effective}
 |R^{(m)}_{n,K}|\le{}&\frac{\alpha^n}{n!}
 \left(m!+\frac{nT_m}{n-K\alpha}\right)\notag\\
 &+\frac{T_m\rho^{-K-1}}{(K+1)^n}
          \frac{\Gamma(n,(K+1)\alpha)}{\Gamma(n)}.
\end{align}
If $n\ge\max\{2,4\alpha^2\}$ and
$K=\lfloor n/\alpha-\sqrt n\rfloor$, then $K\ge1$ and
\begin{equation}\label{reflected:gam:eq:optimized}
 |R^{(m)}_{n,K}|\le
 \left(\frac{m!}{\sqrt{2\pi n}}+
       \frac{T_m}{\alpha\sqrt{2\pi}}+T_m e^{2\alpha^2}\right)
 \left(\frac{e\alpha}{n}\right)^n.
\end{equation}
This bound holds for every $m$, but no useful uniform relative error is
asserted when $m$ grows with $n$.
\end{theorem}
\begin{proof}
The function $f$ decreases from infinity to zero on $(0,1]$:
$\psi'(x)>0$ and $\psi(1)=-\gamma<0$.  Let $x(y)$ denote its inverse,
and set $W_m(x)=\int_0^x B(t)^m\,dt$.  Tonelli's theorem applied to
$f(x)^n=n\int_0^{f(x)}y^{n-1}\,dy$ gives
\begin{equation}\label{reflected:gam:eq:layercake}
 \frac{M_{n,m}}{n!}
 =\frac1{\Gamma(n)}\int_0^\infty y^{n-1}W_m(x(y))\,dy.
\end{equation}

Let $h(q)$ be the local solution of $h=q\phi(h)$ with $h(0)=0$.
Lagrange inversion gives positive coefficients
\[
 h(q)=\sum_{k\ge1}c_kq^k,
 \qquad c_k=\frac1k[u^{k-1}]\phi(u)^k.
\]
The $m=0$ case of Theorem~\ref{reflected:gam:thm:late} proves that its radius
is $\rho$ and that the series converges absolutely on $|q|\le\rho$.
The map $u\mapsto u/\phi(u)$ increases strictly from $0$ to $\rho$
on $[0,\tau]$, so inversion along this interval and continuity imply
$h(\rho)=\tau$.  No other inverse branch is used.

Lagrange--B\"urmann inversion now gives
\[
 W_m(x(y))=\sum_{k\ge1}A_{k,m}e^{-ky}
 \qquad(y\ge\alpha).
\]
Indeed $x(y)=-h(-e^{-y})$ initially for large $y$, and then throughout
$y>\alpha$ by analytic continuation along the strictly monotone real
inverse.  The endpoint follows by absolute convergence.  To prove that
convergence and obtain its quantitative majorant, define
\[
 d_{k,m}=\frac1k[u^{k-1}]\phi(u)^k B(u)^m\ge0.
\]
Coefficientwise absolute values in $B(-u)^m$ are $B(u)^m$.
Consequently $|A_{k,m}|\le d_{k,m}$, while another application of
Lagrange--B\"urmann gives
\[
 \sum_{k\ge1}d_{k,m}\rho^k=W_m(h(\rho))=W_m(\tau)
 \le\tau B(\tau)^m=T_m.
\]
Here $B$ has positive Taylor coefficients and is increasing on $[0,\tau]$.
Thus on $y\ge\alpha$,
\[
 \left|W_m(x(y))-\sum_{k=1}^K A_{k,m}e^{-ky}\right|
 \le T_m\rho^{-K-1}e^{-(K+1)y}.
\]
Its integral in \eqref{reflected:gam:eq:layercake} is the second term of
\eqref{reflected:gam:eq:effective}.

On $0\le y\le\alpha$, use $0\le W_m(x(y))\le M_{m,0}\le m!$.
The last inequality follows from
$0\le\log\Gamma(x)\le-\log x$, since log-convexity gives
$\Gamma(1+x)\le1$ for $0\le x\le1$.  Also, for $k\le K$,
\[
 \int_0^\alpha y^{n-1}e^{-ky}\,dy
 \le\frac{\alpha^n e^{-k\alpha}}{n-k\alpha}.
\]
This follows by differentiating $y^ne^{-ky}$ and using
$n-ky\ge n-k\alpha>0$.  Summing with
$\sum|A_{k,m}|\rho^k\le T_m$ proves the first term of
\eqref{reflected:gam:eq:effective}.

For the stated choice of $K$, $n-K\alpha\ge\alpha\sqrt n$ and
$\theta=(K+1)\alpha/n$ lies in $[1/2,1]$.
The inequality $\theta-1-\log\theta\le2(1-\theta)^2$ gives
\[
 \rho^{-K-1}(K+1)^{-n}
 \le e^{2\alpha^2}(e\alpha/n)^n.
\]
Finally, $\Gamma(n,s)\le\Gamma(n)$ and
$n!\ge\sqrt{2\pi n}(n/e)^n$ yield \eqref{reflected:gam:eq:optimized}.
\end{proof}

\subsection{Sharp truncation for every fixed reflected exponent}
The reflected coefficient expansion and the sharp inverse remainder now
combine. This statement uses the two incoming developments together;
it is stronger than their separate absolute-bound and unreflected results.
Write $\Lambda=-\log\rho$, and retain $\beta_m$ from
\eqref{reflected:gam:eq:beta}.

\begin{corollary}[Half the first omitted reflected term]
\label{reflected:cor:sharp}
Fix an integer $m\ge0$. Let $n,N\to\infty$ through positive integers,
with $\eta=N\Lambda-n$ in a fixed bounded interval. Then
\begin{align}\label{reflected:eq:sharp}
 \frac{M_{n,m}/n!-\sum_{k=1}^{N-1}A_{k,m}/k^n}{A_{N,m}/N^n}
 ={}&\frac12+\frac{3-2\eta}{8N}\notag\\
 &+\frac1{N^2}\left[\frac{\beta_m}{4}
       +\frac{\Lambda(6\eta-13)}{96}\right]+O_m(N^{-3}),
\end{align}
uniformly for bounded $\eta$. In particular the error eventually has
sign $(-1)^{N+m-1}$ and is asymptotic to half the first omitted term.
The denominator is nonzero for every sufficiently large $N$.
\end{corollary}
\begin{proof}
Set $W_m(x)=\int_0^x B(u)^m\,du$ and $G_m(t)=W_m(g(t))$, where $g$ is
the real inverse used in the preceding sharp-moment section.
Lagrange--B\"urmann inversion gives
$G_m(t)=\sum_{k\ge1}A_{k,m}t^k$ near zero. The layer-cake formula
\eqref{reflected:gam:eq:layercake} becomes
\[
 M_{n,m}/n!=\frac1{\Gamma(n)}\int_0^\infty y^{n-1}G_m(e^{-y})\,dy.
\]
Since $W_m$ is holomorphic on $|x|<1$, it can be composed with the
slit-disk inverse of Lemma~\ref{sharp:lem:slit}. At the dominant point
$g(-\rho)=-\tau$ its derivative is
\[
 W_m'(-\tau)=B(-\tau)^m
       =\bigl(\log\Gamma(1+\tau)\bigr)^m=(-b)^m\ne0.
\]
Thus $G_m$ has the same sole dominant square-root singularity, with a
nonzero leading amplitude. Its late coefficients have the relative
coefficient $\beta_m$ established in Theorem~\ref{reflected:gam:thm:late}.
The Cauchy cut-contour argument in Lemma~\ref{sharp:lem:inverse-remainder}
therefore applies to $G_m$ with the same leading kernel
$1/(1+t/\rho)$ and the same $3/(2N)$ endpoint-moment correction.
The sign of the leading jump cancels against $A_{N,m}$ in this ratio.
The second normalized cut coefficient at $t=\rho$ is $\beta_m/4$,
by the same square-root transfer and beta-moment calculation used in
Corollary~\ref{sharp:cor:moment-second}.

The outer real-interval bounds also persist. Above the convergence
radius, the finite coefficient sum is bounded using
$|A_{k,m}|\asymp_m\rho^{-k}k^{-3/2}$, while $0\le G_m(t)\le m!$ for
$0\le t\le1$. The finitely many exceptional early coefficients contribute
only an exponentially small term after normalization. Below the radius,
the absolute Taylor tail gives the corresponding uniform bound.
Average the normalized remainder against the gamma density with shape
$n$ and rate $N$. Its signed moments are exactly those used in the
unreflected calculation. The first two corrections are consequently
\eqref{reflected:eq:sharp}, with $d_1$ replaced by $\beta_m$; the analytic
cut expansion and moments through degree six give the stated uniform
$O_m(N^{-3})$ remainder. Eventual nonvanishing and the sign follow from
\eqref{reflected:gam:eq:late}.
\end{proof}
This fixed-$m$ corollary does not supply uniform relative estimates when
$m$ grows with $n$. That joint regime remains a separate problem.

\subsection{An exact identity coupling all reflected residues}

The coefficients across different $m$ satisfy a finite identity in which
Euler's constant and all odd zeta values cancel.

\begin{proposition}[Reflection sum for the residues]\label{reflected:gam:prop:residue-sum}
For every integer $k\ge1$,
\begin{equation}\label{reflected:gam:eq:residue-sum}
 \sum_{m=0}^{k-1}\frac{k^m}{m!}A_{k,m}
 =\begin{cases}
 0,&k\text{ even},\\[2pt]
 \displaystyle\frac{\pi^{k-1}}{k\,2^{k-1}}
 \binom{k-1}{(k-1)/2},&k\text{ odd}.
 \end{cases}
\end{equation}
\end{proposition}
\begin{proof}
Because $B(x)^m$ has order $m$, extending the sum over $m$ to infinity
does not affect $[x^{k-1}]$.  Euler's reflection formula gives
\[
 \sum_{m=0}^{k-1}\frac{k^m}{m!}A_{k,m}
 =\frac1k[x^{k-1}]\bigl(\Gamma(1+x)\Gamma(1-x)\bigr)^k
 =\frac1k[x^{k-1}]\left(\frac{\pi x}{\sin\pi x}\right)^k.
\]
The coefficient is zero for even $k$.  For odd $k$, compute it as
the residue at $x=0$ of $(\pi/\sin\pi x)^k\,dx$.
The local change of variable $t=\sin\pi x$ gives
\[
 [x^{k-1}]\left(\frac{\pi x}{\sin\pi x}\right)^k
 =\pi^{k-1}[t^{k-1}](1-t^2)^{-1/2},
\]
and the binomial expansion proves \eqref{reflected:gam:eq:residue-sum}.
\end{proof}

For example, the cases $k=2,3$ read
$A_{2,0}+2A_{2,1}=0$ and
$A_{3,0}+3A_{3,1}+\tfrac92 A_{3,2}=\pi^2/6$.
They are exact identities among the asymptotic coefficients, rather
than convergent sums of the moment expansions.  No summation over
the late index $k$ is used.

\subsection{A recorded log-sine conjecture has a formal proof}

There is also a small but useful correction to the logical status of
a claim in the literature.  The author preprint of \cite{gaussian:ACEMKM}
labels a coefficient identity as Conjecture~5.5 and motivates it using
transcendence assumptions on $\log 2$ and $\log\pi$.
For the canonical formal polynomials supplied by the beta integral,
the coefficient identity needs none of those assumptions.  Under the
preprint's assumptions these formal coefficients coincide with its
uniquely specified numerical coefficients, so the result also proves
that stated implication.  We do not assert uniqueness of arbitrary
numerical expressions in zeta values and logarithms, nor priority over
any later treatment.

Define polynomials $\mathcal P_N(X)$ by
\begin{equation}\label{reflected:gam:eq:appell-egf}
 \sum_{N\ge0}\mathcal P_N(X)\frac{z^N}{N!}
 =\exp\left\{Xz+\sum_{j\ge2}
       \frac{(1-2^{1-j})\zeta(j)}jz^j\right\}.
\end{equation}
\begin{proposition}[Unconditional coefficient and translation identities]
\label{reflected:gam:prop:appell}
For $N\ge0$,
\begin{align}
 \int_0^1[-\log\sin(\pi x)]^N\,dx&=\mathcal P_N(\log2),\label{reflected:gam:eq:logsine}\\
 [X^j]\mathcal P_N(X)&=\binom Nj \mathcal P_{N-j}(0),\qquad 0\le j\le N,\notag\\
 \sum_{j=0}^N\binom Nj M_{j,N-j}&=\mathcal P_N(\log(2\pi)).\notag
\end{align}
Moreover, the coefficients admit the finite partition formula
\begin{equation}\label{reflected:gam:eq:appell-partitions}
 \mathcal P_N(X)=N!\!\sum_{r_1+2r_2+\cdots+Nr_N=N}
 \frac{X^{r_1}}{r_1!}
 \prod_{j=2}^N\frac1{r_j!}
 \left(\frac{(1-2^{1-j})\zeta(j)}j\right)^{r_j}.
\end{equation}
\end{proposition}
\begin{proof}
For $\operatorname{Re}z<1$, the beta integral gives
\[
 \int_0^1\sin(\pi x)^{-z}\,dx
 =\frac{\Gamma((1-z)/2)}{\sqrt\pi\,\Gamma(1-z/2)}.
\]
Expanding its logarithm at zero gives the right side of
\eqref{reflected:gam:eq:appell-egf} with $X=\log2$; its degree-$j$ coefficient
for $j\ge2$ is $(1-2^{1-j})\zeta(j)/j$.
Extracting Taylor coefficients proves \eqref{reflected:gam:eq:logsine}.
The factor $e^{Xz}$ gives $[X^j]\mathcal P_N=\binom Nj\mathcal P_{N-j}(0)$ by
formal coefficient extraction.  Euler's reflection formula turns
the sum of mixed moments into the integral of
$(\log\pi-\log\sin\pi x)^N$, which translates $X$ by $\log\pi$.
Expanding each exponential factor separately proves
\eqref{reflected:gam:eq:appell-partitions}.  All formal operations in a fixed
degree involve only finitely many coefficients.
\end{proof}

The constant-column generating function has the additional factorization
\begin{equation}\label{reflected:gam:eq:even-odd-factor}
 \sum_{N\ge0}\mathcal P_N(0)\frac{z^N}{N!}
 =\sqrt{\frac{\tan(\pi z/2)}{\pi z/2}}
 \exp\left\{\sum_{\substack{j\ge3\\j\text{ odd}}}
       \frac{(1-2^{1-j})\zeta(j)}jz^j\right\},
\end{equation}
where the square root is the formal branch with constant term one.
Indeed, multiplying the left side by its value at $-z$ and applying
the reflection formula yields $\tan(\pi z/2)/(\pi z/2)$.
The even logarithmic part is therefore the displayed square root.
This makes the multiplicative factorization of its formal odd-zeta
monomial coefficients explicit, while the even-zeta contribution has
an elementary trigonometric generating function.

\paragraph{Verification and remaining questions.}
The accompanying script \src{code/verify_reflected_moments.py} checks the
three coefficient formulas by exact symbolic arithmetic, checks finite
reflection sums and the formal Appell formulas, and compares independent
quadrature with the fixed-order expansion.  It also compares the late
coefficients with \eqref{reflected:gam:eq:late} and \eqref{reflected:gam:eq:beta}.
The numerical calculations are diagnostics; the preceding analytic and
formal arguments are the proofs.
The fixed-reflected-exponent signed remainder is now supplied by
Corollary~\ref{reflected:cor:sharp}. Uniformity when $m$ increases with
$n$ remains a substantive next question.
In the former regime the endpoint approximation
$B(x)^m\approx\gamma^m x^m$ need not be uniform, so the fixed-$m$
theorem does not settle the simultaneous large-exponent problem.
